% ══════════════════════════════════════════════════════════════════════════════
%  INTRODUCTION
% ══════════════════════════════════════════════════════════════════════════════

\chapter*{Introduction}
\phantomsection
\addcontentsline{toc}{chapter}{Introduction}


% ──────────────────────────────────────────────────────────────────────────────
\section*{About These Notes}
\phantomsection
\addcontentsline{toc}{section}{About These Notes}

These notes accompany the PHYS3071 Astrophysics course in Hong Kong University of Science and Technology (HKUST).
They are intended as a self-contained supplement to the lectures,
providing both conceptual grounding and mathematical rigour at a level
appropriate for readers with a background in special relativity, basic linear algebra, and basic calculus. \\

%\medskip\noindent
The material covers different cosmological topics, including but not limited to the Friedmann equations, and the
energy budget of the $\Lambda$CDM model.
Each topic is structured to serve different learning objectives---from building
physical intuition to solving examination-style problems and exploring
advanced derivations. \\

%\medskip\noindent
\textit{These notes do not replace lecture slides or the recommended textbooks.
They are best used alongside them as a structured study companion.} \\

%\medskip\noindent
If you find any errors or have suggestions for improvement, please feel free to contact me at \texttt{\NoteAuthorEmail}. You can check the latest version of the notes on GitHub: \url{\NoteGitHubRepo}.

% ──────────────────────────────────────────────────────────────────────────────
\section*{Document Conventions}
\phantomsection
\addcontentsline{toc}{section}{Document Conventions}

\subsection*{Block Components}

Each topic is presented using a consistent
\textbf{5-component block structure} designed to support multiple learning styles:

\bigskip

\noindent
\begin{minipage}{0.9\textwidth}
  \begin{tcolorbox}[breakable, colback=GoldBG, colframe=StarGold,
                     leftrule=4pt, arc=0pt,
                     before upper={\textcolor{StarGold}{\textbf{\scshape Physics Intuition}}\par\smallskip}]
    \textbf{Purpose:} Builds conceptual understanding through plain language explanations.

    This section provides the \emph{motivation} behind the physics concept,
    explores its \emph{historical development}, and offers \emph{intuitive explanations}
    that help you grasp what the theory is really about before diving into equations.

    \smallskip
    \textit{Use this when:} You're first encountering a new concept or need to refresh
    the ``big picture'' before tackling mathematics.
  \end{tcolorbox}
\end{minipage}

\bigskip

\noindent
\begin{minipage}{0.9\textwidth}
  \begin{tcolorbox}[colback=TealBG, colframe=NebulaTeal,
                     leftrule=4pt, arc=0pt,
                     before upper={\textcolor{NebulaTeal}{\textbf{\scshape Theory}}\par\smallskip}]
    \textbf{Purpose:} Provides rigorous mathematical treatment of the concept.

    Here you'll find detailed \emph{derivations}, \emph{fundamental assumptions},
    key \emph{equations}, and connections between \emph{mathematics and physical intuition}.
    This is the technical foundation required for exams and deeper problem-solving.

    \smallskip
    \textit{Use this when:} You need to understand how to apply formulas, prepare for exams,
    or work through problems rigorously.
  \end{tcolorbox}
\end{minipage}

\bigskip

\noindent
\begin{minipage}{0.9\textwidth}
  \begin{tcolorbox}[breakable, colback=GreenBG, colframe=ForestGreen,
                     arc=3pt,
                     before upper={\textcolor{ForestGreen}{\textbf{\scshape Question Examples}}\par\smallskip}]
    \textbf{Purpose:} Demonstrates practical application through worked problems.

    Each question example mimics exam-style problems and includes complete
    \emph{solution walkthroughs}. These help you bridge the gap between theory
    and problem-solving, and expose common pitfalls.

    \smallskip
    \textit{Use this when:} Practicing problem-solving, preparing for assessments,
    or checking if you can apply the formulas correctly.
  \end{tcolorbox}
\end{minipage}

\bigskip

\noindent
\begin{minipage}{0.9\textwidth}
  \begin{tcolorbox}[breakable, colback=PurpleBG, colframe=AmethystPurple,
                     leftrule=4pt, arc=0pt,
                     before upper={\textcolor{AmethystPurple}{\textbf{\scshape Foundation}}\par\smallskip}]
    \textbf{Purpose:} Supplies background knowledge and detailed proof steps
    that underpin the main content.

    This section presents \emph{prerequisite mathematics or physics},
    \emph{step-by-step derivations} that are too detailed to include in the Theory block,
    and \emph{intermediate results} needed for a thorough understanding.
    Reading this is optional for the exam but essential for genuine mastery.

    \smallskip
    \textit{Use this when:} A derivation in Theory feels too abrupt, you want to verify
    an intermediate step, or you need to solidify prerequisite concepts.
  \end{tcolorbox}
\end{minipage}

\bigskip

\noindent
\begin{minipage}{0.9\textwidth}
  \begin{tcolorbox}[breakable, colback=GrayBG, colframe=SlateGray,
                     leftrule=2pt, rightrule=2pt, arc=3pt,
                     before upper={\textcolor{SlateGray}{\textbf{\scshape Extra}}\par\smallskip}]
    \textbf{Purpose:} Provides supplementary material that extends \emph{beyond}
    the course syllabus.

    This section offers \emph{interesting asides}, \emph{historical context},
    \emph{connections to current research}, and citations to deeper references.
    The content here is intentionally out-of-scope for the course examination
    but enriches the broader picture of each topic.

    \smallskip
    \textit{Use this when:} You're curious to explore beyond the syllabus, or when a topic
    connects to your personal research interests.
  \end{tcolorbox}
\end{minipage}

\newpage

\noindent\textbf{Suggested Reading Paths:}
\begin{enumerate}
  \item \textbf{Exam Preparation:} \textit{Physics Intuition} $\to$
        \textit{Theory} $\to$ \textit{Question Examples}
  \item \textbf{Deep Learning:} \textit{Physics Intuition} $\to$ \textit{Foundation}
        $\to$ \textit{Theory} $\to$ \textit{Question Examples} $\to$ \textit{Extra}
  \item \textbf{Quick Review:} \textit{Theory} $\to$ \textit{Question Examples}
\end{enumerate}

% ──────────────────────────────────────────────────────────────────────────────
\bigskip \bigskip
\subsection*{Special Annotations}

In addition to the five block components above, two inline annotations appear
throughout the text to flag specific types of information:

\bigskip

\noindent
\begin{minipage}{0.9\textwidth}
  \begin{importantnote}
    \textbf{Purpose:} Alerts you to a result, condition, or sign convention that is
    easy to get wrong or is frequently tested.

    An Important Note signals a \emph{common misconception}, a \emph{critical sign},
    a \emph{non-obvious condition} (e.g.\ a constraint that must hold for an
    equation to be valid), or an \emph{exam trap} worth memorising explicitly.

    \smallskip
    \textit{Treat these as mandatory reading} regardless of which reading path you follow.
  \end{importantnote}
\end{minipage}

\bigskip

\noindent
\begin{minipage}{0.9\textwidth}
  \begin{remark}
    \textbf{Purpose:} Highlights a useful observation, connection, or naming convention.

    A Remark draws attention to a \emph{conceptual link} between results, an
    \emph{elegant reformulation}, a \emph{naming convention} (e.g.\ why a quantity
    carries its particular name), or a \emph{minor but clarifying point} that does
    not rise to the level of an Important Note.

    \smallskip
    \textit{Use this when:} Looking for a deeper understanding of why an equation
    takes the form it does, or for useful mnemonics.
  \end{remark}
\end{minipage}


% ──────────────────────────────────────────────────────────────────────────────
\clearpage
%\bigskip\bigskip
\subsection*{Mathematical Conventions}

% The following conventions are used throughout these notes:

% \medskip
% \noindent\textbf{Metric signature.}\quad
% We adopt the \emph{mostly-plus} (or ``$(-, +, +, +)$'') signature:
% \[
%   g_{\mu\nu} = \operatorname{diag}(-1,\,+1,\,+1,\,+1).
% \]

% \begin{importantnote}
%   Many general-relativity textbooks use the opposite, \emph{mostly-minus}
%   $(+,-,-,-)$ signature. The two conventions differ by an overall sign in
%   the Ricci scalar and stress-energy trace.
%   Always check the signature convention when comparing formulae across sources.
% \end{importantnote}

% \bigskip
% \noindent\textbf{Natural units.}\quad
% In sections on particle physics or early-universe thermodynamics,
% we may set $c = \hbar = k_B = 1$.
% Quantities are then expressed in powers of $\mathrm{eV}$ (or $\mathrm{GeV}$).
% Explicit $c$ and $\hbar$ factors are restored when numerical values are quoted
% in SI units.

% \bigskip
% \noindent\textbf{Key symbols.}
% \begin{center}
% \begin{tabular}{@{}lll@{}}
%   \toprule
%   Symbol & Meaning & Typical units \\
%   \midrule
%   $H_0$                    & Hubble constant today             & $\mathrm{km\,s^{-1}\,Mpc^{-1}}$ \\
%   $h$                      & Dimensionless Hubble: $H_0=100h$  & --- \\
%   $\Omega_i$               & Density parameter for species $i$ & --- \\
%   $a(t)$                   & Scale factor ($a_0 = 1$ today)    & --- \\
%   $z$                      & Cosmological redshift             & --- \\
%   $\rho_\mathrm{crit}$     & Critical density                  & $\mathrm{kg\,m^{-3}}$ \\
%   \bottomrule
% \end{tabular}
% \end{center}

% \begin{remark}
%   The convention $a_0 \equiv a(t_0) = 1$ is standard and simplifies many formulae,
%   e.g.\ $1+z = 1/a(t_e)$.
%   Some older texts normalise $a$ differently; always check the definition
%   when reading the literature.
% \end{remark}

\bigskip \noindent
\textbf{differential notation.}\quad
We use $\dot{x} \equiv \frac{dx}{dt}$ for time derivatives of some variables $x(t)$.

\clearpage
